\documentclass[11pt]{amsart}
\usepackage{amsmath,amssymb,amsthm}
\usepackage[margin=1.15in]{geometry}
\usepackage[hidelinks]{hyperref}
\newtheorem{theorem}{Theorem}[section]
\newtheorem{proposition}[theorem]{Proposition}
\newtheorem{lemma}[theorem]{Lemma}
\theoremstyle{remark}
\newtheorem{remark}[theorem]{Remark}
\newcommand{\R}{\mathbb{R}}
\newcommand{\Z}{\mathbb{Z}}
\newcommand{\E}{\mathbb{E}}
\newcommand{\Prob}{\mathbb{P}}
\newcommand{\Hone}{\mathcal{H}^1}
\newcommand{\spn}{\operatorname{span}}
\newcommand{\doi}[1]{\href{https://doi.org/#1}{doi:\nolinkurl{#1}}}
\newcommand{\arxiv}[1]{\href{https://arxiv.org/abs/#1}{arXiv:\nolinkurl{#1}}}
\setlength{\emergencystretch}{2.5em}

\title[Span length in the Poisson line SIRSN]{The length of the network spanning $k$ random points in the
Poisson line SIRSN}
\author{Yitzchak Shmalo}
\address{Einstein Institute of Mathematics, The Hebrew University of Jerusalem, Jerusalem, Israel}
\email{yspennstate@gmail.com}
\date{November 2026}
\subjclass[2020]{60D05, 60G55, 90B15}
\keywords{scale-invariant random spatial network, SIRSN, improper Poisson line process, random metric, network
length}

\raggedbottom
\begin{document}
\begin{abstract}
Place $k$ independent uniform points in a square of area $k$ and join every pair by its route in a scale-invariant
random spatial network. Aldous asked for a proof that the expected total length of the union of these routes is
asymptotic to $\ell k$, where $\ell$ is the mean length per unit area of the network on a unit-rate Poisson process. We
prove this for the network of geodesics of the planar improper Poisson line process whose line speeds have a power-law
intensity with exponent $\gamma>2$. The lower bound holds for every scale-invariant random spatial network. For the upper bound we show that all
geodesics between points of the unit square lie in the disc of radius $R$ about its centre outside an event of
probability $O(R^{-a})$, for every $0<a<\gamma-1$. Since $\gamma-1>1$, this and a moment bound of Kahn show that for each
$\delta>0$ the expected length of the union at distance more than $\delta\sqrt k$ from the square is $O(\sqrt k)$.
\end{abstract}
\maketitle

\section{Introduction}\label{sec:intro}

The result below was proved by AI agents (Claude, Anthropic). It answers a question of Aldous~\cite{Aldous} for the
network of geodesics of the improper Poisson line process, which Aldous proposed~\cite[Section~4.1]{Aldous}, Kendall
constructed~\cite{Kendall} and Kahn proved to be a scale-invariant random spatial network~\cite{Kahn}; the proof rests
on Theorems~3.1 and~6.1 of Kahn and on the argument that proves his Theorem~5.1.

A scale-invariant random spatial network (SIRSN) assigns to every pair of points $z,z'$ of the plane a random route
$\mathcal R(z,z')$ between them~\cite[Section~2.2]{Aldous}; the notion was introduced by Aldous and
Ganesan~\cite{AG}. The finite-dimensional distributions of the routes are
invariant under translations, rotations and Euclidean scaling, and the route between two points at unit distance has
finite mean length~\cite[(15)]{Aldous}. For a finite set of points, the union
$\spn(z_1,\dots,z_k)=\bigcup_{i,j}\mathcal R(z_i,z_j)$ is called their span. If $\Xi(\lambda)$ is a Poisson process of
intensity $\lambda$ independent of the routes, the span $S(\lambda)$ of its points is a translation-invariant random
set, and its mean length per unit area at $\lambda=1$,
\[
\ell:=\text{mean length per unit area of } S(1),
\]
is required to be finite~\cite[(16)]{Aldous}. The last requirement~\cite[(17)--(20)]{Aldous} concerns the
routes with the unit discs about their endpoints removed: the mean length per unit area of the union of these pieces,
over all pairs of points of $\Xi(\lambda)$, must stay bounded as $\lambda\to\infty$. By~\cite[Lemma~20]{Aldous} it
implies that $\ell$ is finite, so the finiteness of $\ell$ is not a separate requirement. By scaling, $S(\lambda)$ has mean length per unit area
$\ell\lambda^{1/2}$~\cite[(12)]{Aldous}.

Aldous~\cite[Section~8.4.3, Open Problem~35]{Aldous} asks for a proof of the following claim, which he describes
as intuitively obvious but hard to prove: if $Z_1,\dots,Z_k$ are independent and uniform in a square of area $k$, then
\begin{equation}\label{eq:aldous}
\E\,\Hone\bigl(\spn(Z_1,\dots,Z_k)\bigr)\sim\ell k\qquad(k\to\infty),
\end{equation}
that is, the ratio of the two sides tends to $1$. Here $\Hone$ is length (one-dimensional Hausdorff measure), so a
segment shared by several routes is counted once. Aldous locates the difficulty in showing that the span does not,
with small probability, have very large length far from the square.

The first SIRSN constructed with proof is the binary hierarchy model of
Aldous~\cite[Section~3]{Aldous}. The second is the network of geodesics of
the improper Poisson line process, proposed by Aldous~\cite[Section~4.1]{Aldous}; Poisson line patterns had been used
by Aldous and Kendall~\cite{AK} to design networks of near-minimal length on prescribed points. Kendall~\cite{Kendall} constructed the random metric and proved
its basic properties in the plane, and Kahn~\cite{Kahn} proved that its geodesics form a SIRSN in every dimension.
Blanc~\cite{Blanc} computes Hausdorff dimensions for this metric, and Blanc, Curien and Kahn~\cite{BCK} study the
structure of its geodesics in the plane. Kendall~\cite{KendallRRF} studies scale-invariant random flights in the
environment of these lines, and calls the network the Poisson line SIRSN. Each line of the process carries a speed; the speeds have a power-law
intensity whose exponent $\gamma$ we call the speed exponent (Section~\ref{sec:model}), and a geodesic is a path of
least time that moves along the lines at no more than their speeds. We prove~\eqref{eq:aldous} for the Poisson line
SIRSN.

\begin{theorem}\label{thm:main}
Let $\gamma>2$ and let the routes be the geodesics of the planar improper Poisson line process $\Pi$ with speed exponent
$\gamma$. Let $Q_k$ be the square of side $\sqrt k$ centred at $0$, let $Z_1,\dots,Z_k$ be independent, uniform in
$Q_k$ and independent of $\Pi$, and let $L_k=\Hone(\spn(Z_1,\dots,Z_k))$. Then
\[
\lim_{k\to\infty}\frac1k\,\E L_k=\ell .
\]
\end{theorem}

Write $B(c,r)$ for the open disc with centre $c$ and radius $r$, $T(x,y)$ for the least time of a path from $x$
to $y$, $Q=[-\tfrac12,\tfrac12]^2$ for the unit square centred at $0$, and $T^*=\sup_{x,y\in Q}T(x,y)$ for the time diameter of $Q$.
Three results of Kahn~\cite{Kahn} are used: his Theorems~3.1 and~6.1, and the argument that proves his
Theorem~5.1. The time diameter satisfies
$\Prob(T^*>t)\le C_0\exp(-c_0t^{\gamma-1})$, with $C_0,c_0$ depending on $\gamma$~\cite[Theorem~3.1]{Kahn}. For a countable set $\Xi$ independent of the
lines, the total length inside $B(c,r)$ of the geodesics between points of $\Xi$ outside $B(c,2r)$ has $p$th moment at
most $K_pr^p$ for every $p\ge1$, with $K_p$ not depending on $\Xi$~\cite[Theorem~6.1 and its proof]{Kahn}; this is
Lemma~\ref{lem:far} below. The geodesic between two fixed points has finite mean Euclidean
length~\cite[Theorem~5.1]{Kahn}, and the proof shows that this length has finite moments of every order below
$\gamma-1$; so for a fixed pair $x,y$ the geodesic leaves $B(x,R)$ with probability at most $CR^{-a}$, for every
$a<\gamma-1$, with $C$ depending on $a$, $\gamma$ and $|x-y|$.

Fix $\delta>0$. Lemma~\ref{lem:far} alone does not bound the span far from $Q_k$. The points at distance between $2^j\delta\sqrt k$ and
$2^{j+1}\delta\sqrt k$ from $Q_k$ are covered by a bounded number of discs of radius $2^j\delta\sqrt k/4$, so
Lemma~\ref{lem:far} gives a multiple of $2^j\delta\sqrt k$ for the expected length of the span there, and the sum over $j\ge0$
diverges. The bound for a fixed pair has the decay that is missing, but for one geodesic: in Kahn's proof the discs are
centred at the endpoint $x$ and the time bound is $T(x,y)$, and both change with the pair, while the span is a union
of $\binom k2$ geodesics.

We centre the discs at $0$ and replace $T(x,y)$ by $T^*$. If a geodesic between two points of $Q$ reaches distance $R$
from $0$, then for every $\rho\in[3,R]$ it covers Euclidean length at least $\rho-1$ inside $B(0,\rho)$ in time at most
$T^*$, so some line of speed at least $\rho/(2T^*)$ hits $B(0,\rho)$. The computation in Kahn's proof bounds the
probability that this happens at every dyadic radius up to $R$. For $0<a<\gamma-1$ it gives a constant $C_a$,
depending on $a$ and $\gamma$, and for every $R>0$ an event $E'(R)$ of the line process with
\begin{equation}\label{eq:exit-intro}
\Prob(E'(R))\le C_aR^{-a},
\end{equation}
outside which every geodesic between two points of $Q$ lies in $B(0,R)$ (Lemma~\ref{lem:exit}). The rate is that of the bound for a fixed pair; one event
now serves all pairs of points of $Q$.

By scaling, the span of $Z_1,\dots,Z_k$ reaches distance $2^j\delta\sqrt k$ from
$Q_k$ only on an event of probability at most $C_a(2^j\delta)^{-a}$. By H\"older's inequality and Lemma~\ref{lem:far}, its
expected length at distance between $2^j\delta\sqrt k$ and $2^{j+1}\delta\sqrt k$ from $Q_k$ is at most
$C\sqrt k\,2^{j(1-a(1-1/p))}$, where $C$ depends on $\gamma$, $\delta$, $a$ and $p$. Since $\gamma-1>1$ we may take
$a>1$ and then $p$ with $a(1-1/p)>1$, and the sum over $j\ge0$ is at most a multiple of $\sqrt k$. Within distance $\delta\sqrt k$
of $Q_k$ we fix $\varepsilon>0$ and compare with a Poisson process of intensity $1+\varepsilon$: with probability tending to $1$ the span is
contained in the span $S(1+\varepsilon)$, whose mean length per unit area is $\ell(1+\varepsilon)^{1/2}$. This gives
$\limsup_kk^{-1}\E L_k\le\ell$. The matching lower bound (Proposition~\ref{prop:lower}) holds for every SIRSN and
uses only translation invariance, scaling and the monotonicity of the span in the set of points.

Section~\ref{sec:model} recalls the line process and Kahn's results, Section~\ref{sec:lower} proves the lower bound,
Section~\ref{sec:exit} proves Lemma~\ref{lem:exit}, and Section~\ref{sec:upper} completes the proof.

\section{The line process and its geodesics}\label{sec:model}

We follow Kahn~\cite[Section~2]{Kahn}. Let $\mu$ be the
motion-invariant measure on the set of lines of $\R^2$, normalised so that the lines hitting a compact convex set $W$
have measure half the perimeter of $W$; in particular the lines hitting a disc of radius $r$ have measure $\pi r$. For
$\gamma>2$, let $\Pi$ be a Poisson point process of marked lines $(l,v)$ with intensity
$\mu(dl)\,(\gamma-1)v^{-\gamma}\,dv$ on lines times $(0,\infty)$; $v$ is the speed of $l$, and $\gamma$ is the speed
exponent. This $\Pi$ is the improper Poisson line process of Kendall~\cite{Kendall} and Kahn~\cite{Kahn}. The number of
lines of speed at least $v$ hitting a disc of radius $r$ is Poisson with mean
\begin{equation}\label{eq:count}
\pi r\,v^{-(\gamma-1)},
\end{equation}
and the restrictions of $\Pi$ to disjoint sets of lines are independent. A $\Pi$-path is a Lipschitz path
$\xi:[0,T]\to\R^2$ such that for almost every $t$ either $\xi'(t)=0$ or $\xi$ moves along a line of $\Pi$ whose speed
is at least $|\xi'(t)|$; $T$ is its time. Kendall proved~\cite{Kendall} (see~\cite[Theorem~2.4]{Kahn}) that almost
surely every two points $x,y$ are joined by $\Pi$-paths of finite time, that the infimum $T(x,y)$ of their times is
attained, and that finite-time paths have finite Euclidean length. A path attaining $T(x,y)$ is a geodesic. For fixed
$x,y$ the geodesic is almost surely unique~\cite[Theorem~4.4]{Kendall}; by Fubini, if $\Xi$ is a countable random set
independent of $\Pi$, then almost surely all geodesics between points of $\Xi$ are unique, and we write $g_{xy}$ for
them. This network of geodesics is a SIRSN~\cite[Theorem~6.2]{Kahn}; the routes are $\mathcal R(x,y)=g_{xy}$.

\emph{Scaling.} For $s>0$, the image of $\Pi$ under $x\mapsto sx$, with every speed multiplied by $s^{1/(\gamma-1)}$,
has the law of $\Pi$, by~\eqref{eq:count}. Times are multiplied by $s^{(\gamma-2)/(\gamma-1)}$, so geodesics are mapped
to geodesics. The same holds for translations and rotations, without changing speeds.

\emph{Speeds near a disc.} For a disc $B$ let $V(B)$ be the largest speed of a line of $\Pi$ hitting $B$; it is finite
almost surely by~\eqref{eq:count}. A $\Pi$-path that stays in $B$ moves only along lines hitting $B$, so if it covers
Euclidean length $r$ inside $B$ its time is at least $r/V(B)$~\cite[(15)]{Kahn}.

\emph{Time diameter.} With $Q$ and $T^*$ as in Section~\ref{sec:intro}, by~\cite[Theorem~3.1]{Kahn},
applied to the set $Q$ and the lines hitting a neighbourhood of $Q$ (restricting the lines can only increase times),
there are constants $C_0,c_0>0$ depending on $\gamma$ only such that
\begin{equation}\label{eq:diam}
\Prob(T^*>t)\le C_0\exp(-c_0t^{\gamma-1})\qquad(t>0).
\end{equation}

\emph{Geodesics between distant points.} For a countable set $\Xi$, $c\in\R^2$ and $r>0$ let
\[
\Lambda_{c,r}(\Xi)=\Hone\Bigl(\,\bigcup_{x,y\in\Xi\setminus B(c,2r)}g_{xy}\cap B(c,r)\Bigr).
\]

\begin{lemma}\label{lem:far}
Let $\gamma>2$. For every $p\ge1$ there is $K_p<\infty$, depending on $p$ and $\gamma$ only, such that
$\E[\Lambda_{c,r}(\Xi)^p]\le K_p\,r^p$ for every countable random set $\Xi$ independent of $\Pi$, every $c\in\R^2$
and every $r>0$.
\end{lemma}

\begin{proof}
Kahn~\cite[Theorem~6.1]{Kahn} considers the dense countable set $\Xi_\infty$ formed by an increasing sequence of
Poisson processes of intensities $n\in\mathbb N$, and the length inside $B(0,\frac13)$ of the union of the geodesics
between its points with the unit discs about both endpoints removed. He proves that there are
$\eta_{\max},C,C_1>0$ depending on $\gamma$ only such that this length exceeds $C(\ln(C_1/\eta))^2$ with probability
at most $\eta$, for every $\eta<\eta_{\max}$. The first step of his proof bounds that length by
$\Lambda_{0,1/3}(\Xi_\infty)$, and the rest of the proof bounds $\Lambda_{0,1/3}(\Xi_\infty)$ on an event determined by
$\Pi$, using the point set only through the uniqueness of the geodesics between its points. That uniqueness holds for
every countable set independent of $\Pi$, so
$\Prob\bigl(\Lambda_{0,1/3}(\Xi)>C(\ln(C_1/\eta))^2\bigr)\le\eta$ for $\eta<\eta_{\max}$, with the same constants.
Hence $\Prob(\Lambda_{0,1/3}(\Xi)>t)\le C_1\exp(-(t/C)^{1/2})$ for large $t$, and the moments of $\Lambda_{0,1/3}(\Xi)$
are bounded independently of $\Xi$. For general $c$ and $r$, the map $x\mapsto(x-c)/(3r)$ with the speed change above
carries $\Pi$ to a process with the law of $\Pi$, geodesics to geodesics, and $\Xi$ to a countable set $\Xi'$
independent of the new lines, so $\Lambda_{c,r}(\Xi)$ has the law of $3r\,\Lambda_{0,1/3}(\Xi')$.
\end{proof}

\section{The lower bound}\label{sec:lower}

\begin{proposition}\label{prop:lower}
For every SIRSN, with $Q_k$, $Z_1,\dots,Z_k$ and $L_k$ as in Theorem~\ref{thm:main},
\[
\liminf_{k\to\infty}\frac1k\,\E L_k\ge\ell .
\]
\end{proposition}

\begin{proof}
Fix $\varepsilon\in(0,1)$, $\Delta\ge1$ and $M>0$. Let $Z_1,Z_2,\dots$ be independent and uniform in $Q_k$, let $N$ be an
independent Poisson variable with mean $(1-\varepsilon)k$, and let $\Xi$ consist of $Z_1,\dots,Z_N$ together with an
independent Poisson process of intensity $1-\varepsilon$ on $\R^2\setminus Q_k$, all independent of the routes. Then
$\Xi$ is a Poisson process of intensity $\lambda=1-\varepsilon$ and $\Xi\cap Q_k=\{Z_1,\dots,Z_N\}$. The span is
monotone in the point set, so $L_k\ge\Hone(\spn(\Xi\cap Q_k))$ on $\{N\le k\}$.

Call the half-open unit cell $C_z=z+[0,1)^2$, $z\in\Z^2$, a bulk cell if the disc $B_z=B(c_z,\Delta)$ about its centre
$c_z$ lies in $Q_k$. For $\sqrt k\ge2\Delta+2$ there are $n_k\ge(\sqrt k-2\Delta-2)^2$ bulk cells, and $n_k\le k$ because the bulk cells are disjoint
subsets of $Q_k$ of area $1$. For a bulk cell, $\spn(\Xi\cap Q_k)\supset\spn(\Xi\cap B_z)$. The cells are disjoint,
so with $\mathbf 1$ denoting an indicator,
\[
L_k\,\mathbf 1_{\{N\le k\}}\ \ge\ \mathbf 1_{\{N\le k\}}\sum_{z\ \mathrm{bulk}}Y_z,\qquad
Y_z=\min\bigl\{\Hone\bigl(\spn(\Xi\cap B_z)\cap C_z\bigr),M\bigr\}.
\]
By translation invariance of the routes and of $\Xi$, $h_{M,\Delta}:=\E Y_z$ does not depend on $z$. Since $0\le Y_z\le
M$ and $n_k\le k$,
\[
\E L_k\ \ge\ n_kh_{M,\Delta}-Mk\,\Prob(N>k).
\]
By the Chernoff bound $\Prob(N>k)$ decays exponentially in $k$, so $\liminf_kk^{-1}\E L_k\ge h_{M,\Delta}$.

Now let $\Delta\to\infty$. Every route between two points of $\Xi$ lies in $\spn(\Xi\cap B(c_z,\Delta))$ for $\Delta$ large, so these
sets increase to $S(\lambda)$, and $\Hone(\,\cdot\cap C_z)$ increases to $\Hone(S(\lambda)\cap C_z)$ by continuity
of measure from below. By monotone convergence $h_{M,\Delta}\uparrow\E\min\{\Hone(S(\lambda)\cap C_z),M\}$, and as $M\to\infty$
this increases to $\E\,\Hone(S(\lambda)\cap C_z)=\ell\lambda^{1/2}$, the mean length per unit area of $S(\lambda)$
times the area of the cell. Finally let $\varepsilon\to0$.
\end{proof}

\section{Geodesics between points of the unit square}\label{sec:exit}

\begin{lemma}\label{lem:exit}
Let $\gamma>2$ and, for $R>0$, let $E(R)$ be the set of configurations of $\Pi$ for which some geodesic between two
points of $Q$ is not contained in $B(0,R)$. For every $0<a<\gamma-1$ there is $C_a\ge1$, depending on $a$ and $\gamma$, and for every $R>0$ an event
$E'(R)\supset E(R)$ determined by $\Pi$, such that $\Prob(E'(R))\le C_aR^{-a}$.
\end{lemma}

The argument follows Kahn's proof of~\cite[Theorem~5.1]{Kahn}, displays (15)--(22) there, with three changes.
The discs are centred at the centre of $Q$ and not at an endpoint, which halves the speed thresholds. The time
$T(x,y)$ of one geodesic is replaced by $T^*$, so that the bound holds for all pairs of points of $Q$ at once. Kahn
conditions on $\{T(x,y)\le T_n\}$; here the events are determined by the lines alone, and their independence takes
the place of the conditioning. The computation is written out below with his letters.

\begin{proof}[Proof of Lemma~\ref{lem:exit}]
Let $R\ge3$ and suppose $E(R)$ occurs: $g$ is a geodesic from $x\in Q$ to $y\in Q$ through a point $w$ with $|w|\ge
R$. Let $3\le\rho\le R$ and $r=\rho-1$. Since $|x|<1$, $|w-x|>R-1\ge r$, so the part of $g$ from $x$ to its first exit
from $B(x,r)$ has Euclidean length at least $r$ and lies in $B(x,r)\subset B(0,\rho)$. Its time is at most $T(x,y)\le
T^*$, so $V(B(0,\rho))\ge r/T^*\ge\rho/(2T^*)$. Hence
\begin{equation}\label{eq:constraint}
E(R)\cap\{T^*\le\tau\}\subset\bigcap_{3\le\rho\le R}\bigl\{V(B(0,\rho))\ge\rho/(2\tau)\bigr\}
\end{equation}
for every $\tau>0$; the value of $\tau$ is chosen below.

Fix $r_0\ge4$ and an integer $m\ge0$ with $2^mr_0\le R$. Write $r_l=2^lr_0$, $V_l=V(B(0,r_l))$ and
$v_l=r_l/(2\tau)$, so $v_l=2^lv_0$, and let $G=\bigcap_{l=0}^m\{V_l\ge v_l\}$. On $G$ define $l_0=0$ and, as long as
some $l\le m$ has $v_l>V_{l_i}$, let $l_{i+1}$ be the least such $l$; it exceeds $l_i$ because $V_{l_i}\ge v_{l_i}$.
The recursion stops at some $l_I$, and we set $l_{I+1}=m+1$. Then $V_{l_i}\ge v_{l_{i+1}-1}$ for $0\le i\le I$, and
$V_{l_i}<v_{l_{i+1}}$ for $i<I$.

For a fixed integer $I\in[0,m]$ and a fixed sequence $0=l_0<l_1<\dots<l_I\le m$, with $l_{I+1}=m+1$, let $D_i$
($0\le i\le I$) be the event that some line of speed at least $v_{l_{i+1}-1}$ hits $B(0,r_{l_i})$ and, if $i\ge1$,
misses $B(0,r_{l_{i-1}})$. If $G$ occurs and the recursion above produces this sequence, then every $D_i$ occurs:
$V_{l_i}\ge v_{l_{i+1}-1}$ gives a line of that speed hitting $B(0,r_{l_i})$, and for $i\ge1$ it misses
$B(0,r_{l_{i-1}})$ because $V_{l_{i-1}}<v_{l_i}\le v_{l_{i+1}-1}$. So $G$ is contained in the union over $I$ and over
the sequences of $D_0\cap\dots\cap D_I$. For a fixed sequence, $D_i$ is determined by the lines hitting
$B(0,r_{l_i})$ but not $B(0,r_{l_{i-1}})$ (all lines hitting $B(0,r_{l_0})$ for $i=0$). These sets of lines are
disjoint, so $D_0,\dots,D_I$ are independent, and by~\eqref{eq:count}
\[
\Prob(D_i)\le\pi r_{l_i}v_{l_{i+1}-1}^{-(\gamma-1)}
=2^{-(\gamma-1)(l_{i+1}-l_i-1)}\pi r_{l_i}^{2-\gamma}(2\tau)^{\gamma-1}
\le2^{-(\gamma-1)(l_{i+1}-l_i)}\beta,
\]
where $\beta=2^{\gamma-1}\pi r_0^{2-\gamma}(2\tau)^{\gamma-1}$, using $r_{l_i}\ge r_0$ and $\gamma>2$. The product of
the factors $2^{-(\gamma-1)(l_{i+1}-l_i)}$ over $0\le i\le I$ is $2^{-(\gamma-1)(m+1)}$. Summing over the sequences, of
which there are $\binom mI$ for each $I$,
\[
\Prob(G)\le\sum_{I=0}^m\binom mI2^{-(\gamma-1)(m+1)}\beta^{I+1}\le\bigl(2^{-(\gamma-1)}(1+\beta)\bigr)^{m+1}.
\]

Now fix $\kappa$ with $a<\kappa<\gamma-1$, and put $\tau=(\theta\ln R)^{1/(\gamma-1)}$ with $\theta=a/c_0$, so that
$\Prob(T^*>\tau)\le C_0R^{-a}$ by~\eqref{eq:diam}. Choose
\[
r_0=\max\Bigl\{4,\Bigl(\frac{2^{\gamma-1}\pi(2\tau)^{\gamma-1}}{2^{\gamma-1-\kappa}-1}\Bigr)^{1/(\gamma-2)}\Bigr\},
\]
so that $1+\beta\le2^{\gamma-1-\kappa}$ and $\Prob(G)\le2^{-\kappa(m+1)}$. Since $r_0=O((\ln R)^{1/(\gamma-2)})$, there
is $R_a\ge4$ such that $r_0\le R$ and $(r_0/R)^\kappa\le R^{-a}$ for all $R\ge R_a$. For such $R$ let $m$ be the
largest integer with $2^mr_0\le R$. Then $2^{m+1}>R/r_0$, the radii $r_0,\dots,r_m$ lie in $[4,R]$, and
by~\eqref{eq:constraint} the event $E'(R)=\{T^*>\tau\}\cup G$ contains $E(R)$ and satisfies
\[
\Prob(E'(R))\le\Prob(T^*>\tau)+\Prob(G)\le C_0R^{-a}+(r_0/R)^{\kappa}\le(C_0+1)R^{-a}.
\]
For $R<R_a$ let $E'(R)$ be the whole probability space. The lemma holds with $C_a=\max\{C_0+1,R_a^{\,a}\}$.
\end{proof}

\section{The upper bound}\label{sec:upper}

\begin{proof}[Proof of Theorem~\ref{thm:main}]
By Proposition~\ref{prop:lower} it remains to show $\limsup_kk^{-1}\E L_k\le\ell$. Fix $\varepsilon,\delta\in(0,1)$.

\emph{Poissonization.} In this proof, let $Z_1,Z_2,\dots$ be independent and uniform in $Q_k$, let $N$ be an independent Poisson
variable with mean $(1+\varepsilon)k$, and let $\Xi$ consist of $Z_1,\dots,Z_N$ together with an independent Poisson
process of intensity $\lambda=1+\varepsilon$ on $\R^2\setminus Q_k$, everything independent of $\Pi$. On $\{N\ge k\}$,
$L_k\le L':=\Hone(\spn(\Xi\cap Q_k))$, and $L_k$ is independent of $N$, so
\[
\E L_k\,\Prob(N\ge k)=\E[L_k;N\ge k]\le\E L'.
\]
Since $\Prob(N\ge k)\to1$, it suffices to bound $\E L'$.

\emph{Near the square.} Let $U_0=\{w:\operatorname{dist}(w,Q_k)<\delta\sqrt k\}$, of area $k(1+4\delta+\pi\delta^2)\le
(1+2\delta)^2k$. Since $\spn(\Xi\cap Q_k)\subset S(\lambda)$, the span of all of $\Xi$,
\[
\E\,\Hone\bigl(\spn(\Xi\cap Q_k)\cap U_0\bigr)\le\ell\lambda^{1/2}(1+2\delta)^2k .
\]

\emph{Far from the square.} For $j\ge0$ let $\rho_j=2^j\delta\sqrt k$ and $A_j=\{w:\rho_j\le\operatorname{dist}(w,Q_k)
<2\rho_j\}$; the sets $U_0,A_0,A_1,\dots$ cover the plane. Fix $0<a<\gamma-1$. Since $0\in Q_k$, a point $w$ of $A_j$
has $|w|\ge\rho_j$, so if $\spn(\Xi\cap Q_k)$ meets $A_j$, some geodesic between two points of $Q_k$ is not contained
in $B(0,\rho_j)$. Scaling by $k^{-1/2}$ maps $Q_k$ to $Q$ and geodesics to geodesics, so this happens only on the
event $F_j$, the image under the scaling of $E'(2^j\delta)$ in Lemma~\ref{lem:exit}, and
$\Prob(F_j)\le q_j:=\min\{1,C_a(2^j\delta)^{-a}\}$.

Let $s_j=\rho_j/4$ and let $\mathcal C_j\subset A_j$ be a maximal set of points at mutual distance at least $s_j$.
The discs $B(c,s_j)$, $c\in\mathcal C_j$, cover $A_j$, and the discs $B(c,s_j/2)$ are disjoint and lie in the square
of side $\sqrt k+5\rho_j$ centred at $0$, so
\[
\#\mathcal C_j\le\frac{(\sqrt k+5\rho_j)^2}{\pi s_j^2/4}=\frac{64}{\pi}\Bigl(\frac1{2^j\delta}+5\Bigr)^2\le
\frac{64}{\pi}\Bigl(\frac1\delta+5\Bigr)^2=:\nu_\delta .
\]
Every point of $Q_k$ is at distance at least $\rho_j=4s_j>2s_j$ from each $c\in\mathcal C_j$, so
$\spn(\Xi\cap Q_k)\cap B(c,s_j)$ is contained in the set whose length is $\Lambda_{c,s_j}(\Xi)$. Hence
\[
\Hone\bigl(\spn(\Xi\cap Q_k)\cap A_j\bigr)\le\mathbf 1_{F_j}\sum_{c\in\mathcal C_j}\Lambda_{c,s_j}(\Xi).
\]
By H\"older's inequality and Lemma~\ref{lem:far}, for $p>1$,
\[
\E\,\Hone\bigl(\spn(\Xi\cap Q_k)\cap A_j\bigr)\le \nu_\delta K_p^{1/p}s_j\,q_j^{1-1/p}
=\frac{\nu_\delta K_p^{1/p}\delta}{4}\sqrt k\;2^j\min\bigl\{1,C_a(2^j\delta)^{-a}\bigr\}^{1-1/p}.
\]
Since $\gamma-1>1$ we may take $a\in(1,\gamma-1)$ and then $p$ so large that $a(1-1/p)>1$. The series over $j\ge0$
then converges, and its sum is at most $C(\delta)\sqrt k$, with $C(\delta)$ independent of $k$.

\emph{Conclusion.} Adding the two parts,
\[
\E L_k\le\frac{\ell(1+\varepsilon)^{1/2}(1+2\delta)^2k+C(\delta)\sqrt k}{\Prob(N\ge k)},
\]
so $\limsup_kk^{-1}\E L_k\le\ell(1+\varepsilon)^{1/2}(1+2\delta)^2$. Letting $\delta,\varepsilon\to0$ completes the
proof.
\end{proof}

\section*{Notation}
\begin{tabular}{@{}p{0.22\textwidth}p{0.52\textwidth}p{0.18\textwidth}@{}}
symbol & meaning & first appears\\
$\mathcal R(z,z')$, $\spn$ & route between $z$ and $z'$; union of the routes between given points & Section~\ref{sec:intro}\\
$S(\lambda)$, $\ell$ & span of a Poisson process of intensity $\lambda$; mean length per unit area of $S(1)$ &
Section~\ref{sec:intro}\\
$\Hone$ & length (one-dimensional Hausdorff measure) & \eqref{eq:aldous}\\
$Q_k$, $L_k$ & square of side $\sqrt k$ centred at $0$; length of the span of $Z_1,\dots,Z_k$ &
Theorem~\ref{thm:main}\\
$\Pi$, $\gamma$ & line process, speed exponent & Section~\ref{sec:intro}\\
$\mu$ & invariant measure on lines & Section~\ref{sec:model}\\
$T(x,y)$, $B(c,r)$ & least time from $x$ to $y$; open disc & Section~\ref{sec:intro}\\
$g_{xy}$, $V(B)$ & the geodesic from $x$ to $y$; largest speed of a line hitting the disc $B$ & Section~\ref{sec:model}\\
$Q$, $T^*$, $C_0$, $c_0$ & unit square centred at $0$; its time diameter; constants in~\eqref{eq:diam} &
Section~\ref{sec:intro}\\
$\Lambda_{c,r}(\Xi)$ & length in $B(c,r)$ of the geodesics between points of $\Xi$ outside $B(c,2r)$ &
Section~\ref{sec:model}\\
$K_p$, $\Xi$ & moment constant of $\Lambda_{c,r}$; a countable or Poisson point set & Section~\ref{sec:intro}\\
$\delta$, $\varepsilon$ & width of the neighbourhood of $Q_k$ in units of $\sqrt k$; intensity correction &
Section~\ref{sec:intro}\\
$N$ & Poisson count used to compare with $Z_1,\dots,Z_k$ & Section~\ref{sec:lower}\\
$C_z$, $B_z$, $\Delta$, $Y_z$, $M$, $h_{M,\Delta}$, $n_k$ & unit cells, discs about their centres, radius, truncated lengths,
truncation level, their mean, number of bulk cells & Section~\ref{sec:lower}\\
$E(R)$ & set of configurations with a geodesic between points of $Q$ leaving $B(0,R)$ & Lemma~\ref{lem:exit}\\
$E'(R)$, $C_a$ & the event of $\Pi$ containing $E(R)$; its constant & Section~\ref{sec:intro}\\
$r_l$, $V_l$, $v_l$, $G$, $D_i$, $\beta$, $\tau$, $\kappa$ & radii, speeds and events of the proof of Lemma~\ref{lem:exit} &
Section~\ref{sec:exit}\\
$U_0$, $A_j$, $\rho_j$, $F_j$, $q_j$ & neighbourhood of $Q_k$, dyadic annuli around it, scaled copies of $E'$, their probabilities & Section~\ref{sec:upper}\\
$s_j$, $\mathcal C_j$, $\nu_\delta$ & radius, centres and number of the discs covering $A_j$ & Section~\ref{sec:upper}\\
\end{tabular}

\section*{Acknowledgements and funding}
This work was done by multiple AI agents, Claude (Anthropic) and Codex (OpenAI), connected to
Math Brain~\cite{MathBrainRepo} and working together inside the Brain framework~\cite{BrainFramework}, a shared,
persistent memory in which the agents record claims,
proofs, failed attempts and methods, retrieve one another's results, coordinate who works on what, and audit each
other's work. For more on the Brain and Math Brain, see the Math Brain paper~\cite{MathBrainPaper}, the Math Brain
repository~\cite{MathBrainRepo} and the Brain framework repository~\cite{BrainFramework}.

The research presented in this paper was supported by the European Research Council (ERC) under the European Union's
Horizon Europe research and innovation programme (grant agreement No.~101041711), by the Simons Foundation as part of
the Collaboration on the Mathematical and Scientific Foundations of Deep Learning, by Heights Labs, by Convex Nexus
Capital, by the Israel Science Foundation (grant number 2258/19), and by the Israel Science Foundation (ISF Grant
4101/25). Computational resources were provided in part by Google Cloud through its education credits programme
(credit 530550103).

\section*{Declaration of competing interest}
The author holds equity in Heights Labs and in Convex Nexus Capital and has a management role at Convex Nexus Capital.
These entities are included in the funding declaration above.

\section*{Data availability}
No data or code were produced for this work.

\section*{Declaration of generative AI and AI-assisted technologies}
Claude agents (Anthropic) proved Theorem~\ref{thm:main} and prepared the manuscript, and other Claude agents refereed
the proofs independently. The problem is due to Aldous~\cite{Aldous}. The network was proposed by
Aldous~\cite{Aldous} and constructed by Kendall~\cite{Kendall}. The proof rests on Theorems~3.1 and~6.1 of
Kahn~\cite{Kahn} and on the argument that proves his Theorem~5.1; these results are theirs. New here are the bound
for this network and Lemma~\ref{lem:exit}, which adapts Kahn's argument to all pairs of points of $Q$ at once. The author directed the work and is responsible for
the content of the paper.

\begin{thebibliography}{99}
\bibitem{Aldous} D.~J.~Aldous, \emph{Scale-invariant random spatial networks}, Electron. J. Probab. \textbf{19}
(2014), 1--41, \doi{10.1214/EJP.v19-2920}; section, display and problem numbers as in \arxiv{1204.0817v1}.
\bibitem{AG} D.~J.~Aldous and K.~Ganesan, \emph{True scale-invariant random spatial networks}, Proc. Natl. Acad. Sci.
USA \textbf{110} (2013), 8782--8785, \doi{10.1073/pnas.1304329110}.
\bibitem{AK} D.~J.~Aldous and W.~S.~Kendall, \emph{Short-length routes in low-cost networks via Poisson line
patterns}, Adv. in Appl. Probab. \textbf{40} (2008), 1--21, \doi{10.1239/aap/1208358883}.
\bibitem{Blanc} G.~Blanc, \emph{Fractal properties of Aldous--Kendall random metric}, preprint, 2022,
\arxiv{2207.03349}.
\bibitem{BCK} G.~Blanc, N.~Curien and J.~Kahn, \emph{Geodesics in planar Poisson roads random metric}, preprint,
2024, \arxiv{2407.07887}.
\bibitem{Kahn} J.~Kahn, \emph{Improper Poisson line process as SIRSN in any dimension}, Ann. Probab. \textbf{44}
(2016), 2694--2725, \doi{10.1214/15-AOP1032}; theorem and display numbers as in \arxiv{1503.03976}.
\bibitem{KendallRRF} W.~S.~Kendall, \emph{Rayleigh random flights on the Poisson line SIRSN}, Electron. J.
Probab. \textbf{25} (2020), 1--36, \doi{10.1214/20-EJP526}.
\bibitem{Kendall} W.~S.~Kendall, \emph{From random lines to metric spaces}, Ann. Probab. \textbf{45} (2017),
469--517, \doi{10.1214/14-AOP935}; theorem numbers as in \arxiv{1403.1156}.
\bibitem{MathBrainRepo} Y.~Shmalo, \emph{Math Brain: public snapshot, code and tests}, 2026,
\url{https://github.com/yspennstate/math-brain-public/tree/8c9a6bd478c4ec4082347151b7ae16db50c917dd}.
\bibitem{MathBrainPaper} Y.~Shmalo, \emph{Math Brain: shared, auditable memory for multi-agent mathematical research},
October 2026, \doi{10.5281/zenodo.24150837},
\url{https://github.com/yspennstate/math-brain-public/tree/main/papers/math-brain-2026-10-24}.
\bibitem{BrainFramework} Y.~Shmalo, \emph{Brain Framework}, public software repository, 2026,
\url{https://github.com/yspennstate/brain-framework-public/tree/5629e644cdd510173ae97e9aa92642dabd22dc57}.
\end{thebibliography}

\end{document}
